\chapter{Evaluation}
\label{ch:eval}

This chapter evaluates the guidance method in two stages. We first compare
different choices for the guidance-strength profile through a structured set of
quantitative experiments and select the configuration used in the remainder of
the thesis. With this configuration fixed, we then investigate how the method
behaves under different scenario specifications, building an empirical
understanding of how the requested target affects the resulting trajectories.

Throughout the experiments, we use a small set of diagnostics to verify that
guidance behaves as intended and to identify anomalous runs before interpreting
their results.

The first evaluation, T1, provides a systematic comparison of candidate
guidance-strength profiles across multiple experimental settings. We assess
each profile using a common set of complementary quantitative criteria and
aggregate the resulting comparisons to identify a configuration that performs
robustly across the considered settings.

In T2, we fix the selected guidance configuration and vary the parameters that
define the scenario of interest. Rather than producing another standardized
ranking, this analysis examines how changes in the specification affect the
generated trajectories. Quantitative measures are used alongside the resulting
fields and trajectories where they help characterize and explain the observed
behavior.

The findings of these evaluations provide the basis for constructing and
analyzing a complete high-temperature storyline in the following chapter.
\section{T0: Experiment diagnostics}
\label{sec:experiment-diagnostics}

Before interpreting an experiment, we inspect a small set of diagnostics to
verify that the guidance procedure behaved as intended. Specifically, we check

\begin{itemize}[noitemsep]
    \item whether each guided member reaches its prescribed target within a
    defined tolerance;
    \item how the targeted regional quantity evolves over the generative
    process, in particular for oscillatory or irregular convergence;
    \item whether the final guided field exhibits conspicuous spatial
    artifacts relative to the unguided field; and
    \item the spatial structure of the isolated guidance effect.
\end{itemize}

For the final diagnostic, we compare the guided state with its matched
unguided-under-guided-context counterpart. Since
$x_n^{\mathrm{gui}}$ and $x_n^{\mathrm{ung}\mid\mathrm{gui}}$ share the same
conditioning states, deterministic forecast, and initial residual noise,

\[
    \Delta x_n^{\mathrm{guidance}}
    :=
    x_n^{\mathrm{gui}}
    -
    x_n^{\mathrm{ung}\mid\mathrm{gui}}
    =
    r_n^{\mathrm{gui}}
    -
    r_n^{\mathrm{ung}\mid\mathrm{gui}},
\]

isolates the effect of guidance within the current residual flow.

These diagnostics provide a first check that an experiment is suitable for the
more detailed evaluations that follow.

\section{T1: Guidance-profile evaluation}
\label{sec:guidance-profile-evaluation}

We compare the candidate guidance-strength profiles through four complementary
properties of the resulting guidance effect. Each metric is aggregated over
ensemble members, guided weather steps, and the experiments included in the
evaluation sweep. The resulting scores are combined in a leaderboard to compare
the different flow-time profiles under a common set of criteria.


\subsection{T1.1: Multivariate propagation of guidance}

The first criterion measures how strongly an intervention targeted at regional
surface temperature propagates to the remaining atmospheric state. We consider
the guidance effect

\[
    \Delta x_n^{\mathrm{guidance}}
    =
    x_n^{\mathrm{gui}}
    -
    x_n^{\mathrm{ung}\mid\mathrm{gui}},
\]

and quantify its magnitude both within the target region and over the complete
spatial domain. The resulting masked and global push scores summarize how
strongly the intervention extends beyond the directly targeted quantity.


\subsection{T1.2: Deviation from the unguided trajectory}

We next measure how far the guided state is displaced from its matched
unguided-under-guided-context counterpart. The comparison is performed in a
low-dimensional representation of the atmospheric state obtained through
principal component analysis. The resulting distance provides a common measure
of how strongly each guidance profile modifies the trajectory relative to the
corresponding unguided generation.


\subsection{T1.3: Spatial deviation from the target mask}

The guidance objective is defined using a simple rectangular spatial mask. A
guided response that reproduces this pattern directly would indicate a strong
imprint of the prescribed objective on the generated field. We therefore
compare the spatial distribution of the guidance effect with the target mask
using total variation distance.

We compute this distance separately for the targeted surface-temperature field
and across the complete atmospheric state. Larger values indicate a stronger
departure from the imposed mask pattern.


\subsection{T1.4: Guidance-effect diversity across ensemble members}

Finally, we assess whether the spatial guidance effect remains dependent on the
stochastic realization or converges toward a common perturbation across
ensemble members. We quantify the spread of the guidance footprint across
members, both for the target field and across all atmospheric fields.

Larger footprint spread indicates greater member-to-member diversity in the
response to guidance.

% TODO:
% - Define the candidate relative guidance-strength profiles.
% - Define the experimental settings over which they are compared.
% - Introduce the quantitative evaluation criteria.
% - Explain how results are aggregated across experiments/members.
% - Rank the candidate profiles.
% - Select the profile used for the remaining experiments.


\section{T2: Scenario-specification analysis}
\label{sec:scenario-specification-analysis}

% TODO:
% - Fix the guidance configuration selected in T1.
% - Vary the scenario-specification parameters.
% - Investigate target intensity / trajectory.
% - Investigate spatial region.
% - Investigate target variable if included in the final experiments.
% - Combine trajectory and field visualizations with quantitative measures
%   where useful for supporting the analysis.
% - Extract practical insights for the complete storyline experiment.